\begin{center}
Problema F

Monitorando Pragas

{\small Nome do arquivo: F.c, F.cpp, F.java, ou F.py}

{\large Timelimit: $1$}

{\tiny Por Emanuel Silva Araujo}

\end{center}

\footnotesize

O agricultor Seu Jayme sofre constantemente com infestações de insetos em sua lavoura de soja. Para reduzir perdas e otimizar o manejo, ele procurou a EMBRAPII para desenvolver uma ferramenta de monitoramento diário de pragas.
Sua tarefa como programador é implementar um programa que leia os valores de população diária e identifique situações de risco, emitindo alertas de acordo com os seguintes critérios:
%Critérios de Alerta
    \begin{itemize}
        \item  Alerta por pico diário [Pico]
        \begin{itemize}
            \item Se em qualquer dia a população registrada ultrapassar 200 insetos.
        \end{itemize}
        \item  Alerta por variação em 7 dias [7 Dias]
        \begin{itemize}
            \item Se a diferença percentual entre o valor mínimo e o máximo em um intervalo de 7 dias consecutivos for maior ou igual a 20\%.
        \end{itemize}
        \item  Alerta por crescimento contínuo [Crescimento]
        \begin{itemize}
            \item Se ocorrerem 4 dias seguidos de crescimento, e em cada dia a população for pelo menos 5\% maior que a do dia anterior.
        \end{itemize}
    \end{itemize}
        
\vspace{0.3cm}
\textbf{Entrada}

Um número inteiro $(0 <= N <= 360)$, representando a quantidade de dias monitorados (entre os dias $0$ e $360$).
Uma sequência de $N$ números inteiros $(0 < X <= 5000)$, cada um representando a estimativa da população de insetos no respectivo dia.

\vspace{0.3cm}
\textbf{Saída}

    Para cada alerta identificado, o programa deve imprimir:
    \begin{itemize}
        \item O tipo do alerta;
        \item O(s) dia(s) correspondente(s);
        \item O valor percentual (quando aplicável), com duas casas decimais.
    \end{itemize}
    Caso não tenha nenhum alerta, deve imprimir uma mensagem de "Nenhum Alerta".
    Para dias que tiveram mais de um alerta a ordem de impressão é: $[Pico] > [7 Dias] > [Crescimento]$.
 
\begin{table}[h]
    \centering
    \footnotesize
    \begin{tabular}{|p{8cm}|p{8cm}|}
        \hline
        \begin{center}
            \vspace{-0.3cm}
            \textbf{Exemplos de Entrada}
            \vspace{-0.3cm}
        \end{center} 
         &  
        \begin{center}
            \vspace{-0.3cm}
            \textbf{Exemplos de Saída}
            \vspace{-0.3cm}
        \end{center} \\ \hline
		    \texttt{10} & \texttt{[Crescimento] 1-4: maior que >= 5\% } \\
                \texttt{50 60 70 80 90 200 210 220 180 170} & \texttt{[Crescimento] 1-5: maior que >= 5\% } \\
& \texttt{[Crescimento] 1-6: maior que >= 5\% } \\
& \texttt{[Pico] 7: 210 insetos} \\
& \texttt{[7 Dias] 1-7: 320.00\%} \\
& \texttt{[Crescimento] 1-7: maior que >= 5\% } \\
& \texttt{[Pico] 8: 220 insetos} \\
& \texttt{[7 Dias] 2-8: 266.67\%} \\
& \texttt{[7 Dias] 3-9: 214.29\%} \\
& \texttt{[7 Dias] 4-10: 175.00\%} \\


         \hline
   		    \texttt{9}   & \texttt{[Crescimento] 1-4: maior que >= 5\% } \\
            \texttt{100 105 112 118 90 96 100 120 80} & \texttt{[7 Dias] 2-8: 33.33\%} \\
         \hline    
   		    \texttt{8}   & \texttt{Nenhum Alerta} \\
            \texttt{118 120 118 140 140 140 130 100} & \\
         \hline    
    \end{tabular}
    \caption{Exemplos de entradas e saídas}

\end{table}

\normalsize

\newpage

\begin{center}
\lstinputlisting{abobora_25/F/F5.cpp}
\end{center}
